% ======================================================================
% CHAPTER 2: LITERATURE REVIEW / ANALYSIS
% ======================================================================
\chapter{Literature Review / Analysis}

\section{Related Work}
Research at the intersection of natural language processing and clinical document understanding has grown substantially over the past decade. Early work focused on named entity recognition in clinical text, with systems such as cTAKES (Clinical Text Analysis and Knowledge Extraction System) demonstrating the feasibility of automated extraction of clinical concepts from unstructured reports. These rule-based and statistical systems, however, required significant domain-specific engineering and lacked the ability to generate natural, explanatory prose. Similarly, the Apache cTAKES project, developed at the Mayo Clinic, provided a robust framework for extracting clinical information from free-text electronic medical records, but remained fundamentally a structured data extraction pipeline rather than a patient-facing interpretation tool.

The advent of transformer-based models fundamentally changed the landscape. BioBERT (Lee et al., 2020) established that pre-training on biomedical corpora significantly improves performance on clinical NLP tasks including relation extraction, question answering, and entity recognition. The authors demonstrated that BioBERT, pre-trained on PubMed abstracts and PMC full-text articles, outperformed BERT by a substantial margin across all evaluated biomedical NLP benchmarks. PubMedBERT (Gu et al., 2021) extended this work by showing that domain-specific pre-training from scratch---rather than continued pre-training of a general-domain BERT---yields superior performance on biomedical tasks. BioMedLM, developed by Stanford CRFM and MosaicML, further extended this line of work to generative settings by pre-training a 2.7-billion-parameter GPT-style model on the full PubMed corpus, enabling the generation of coherent biomedical text from trained parameters alone.

The introduction of Retrieval-Augmented Generation (Lewis et al., 2020) provided a foundational framework for combining the fluency of large language models with the factual grounding of document retrieval. The RAG architecture, as originally proposed, combined a BART-based sequence-to-sequence model with a dense passage retriever (DPR), enabling the model to retrieve and condition on relevant documents during generation. Subsequent work demonstrated that RAG architectures substantially reduce hallucination rates in knowledge-intensive tasks, making them particularly well-suited for medical document analysis where factual accuracy is paramount. Variants such as REPLUG (Shi et al., 2023) enhanced this paradigm by treating the LM as a black-box and applying an iterative retrieval loop that refines the retrieved set based on the LM's own perplexity signal. Self-RAG (Asai et al., 2023) introduced self-reflection mechanisms that allow the model to decide when retrieval is needed and to critique its own generated outputs, further improving both retrieval quality and output faithfulness.

Recent systems such as Med-PaLM 2 (Singhal et al., 2023) demonstrated that a 540-billion-parameter LLM fine-tuned on medical exam questions could achieve performance approaching that of human experts on the USMLE-style benchmarks, scoring 86.5\% on the MedQA dataset. Similarly, GPT-4 demonstrated significant medical knowledge capabilities, scoring in the top 10\% of USMLE Step 2 CK examinees. However, these frontier models operate in a general question-answering paradigm and are not specifically designed for the structured, patient-facing task of parsing an uploaded diagnostic report and producing a formatted analysis with layered content (summary, observations, flagged values, guidance). Consumer-facing medical interpretation tools such as Buoy Health and Ada Health provide symptom-checking and triage guidance through interactive chatbot interfaces, but operate primarily through structured questionnaires rather than document upload and analysis. WebMD and other health information portals offer comprehensive medical encyclopaedia content but do not provide personalised report interpretation. ReportRx targets the specific, underserved application of direct diagnostic report interpretation using a grounded, multi-model pipeline architecture.

\section{Comparative Study}
The following table summarises existing tools relevant to this domain and compares their core capabilities against the proposed ReportRx system:

\begin{table}[h]
\centering
\caption{Comparative Study of Related Systems}
\begin{tabular}{|p{1.9cm}|p{1.9cm}|p{2.2cm}|p{2.4cm}|p{1.6cm}|}
    \hline
    \textbf{Tool / System} & \textbf{Approach} & \textbf{Strength} & \textbf{Limitation} & \textbf{Report Upload?} \\
    \hline
    cTAKES & Rule-based NLP & High precision on structured EHR data & No generative output; requires expert setup & No \\
    BioBERT & Pre-trained transformer & Strong clinical entity extraction & Encoder-only; cannot generate explanations & No \\
    Med-PaLM 2 & Fine-tuned frontier LLM & Near-expert medical QA performance & Proprietary; not designed for report upload & No \\
    General ChatGPT & General LLM (no RAG) & Fluent explanations & No document grounding; hallucination risk & Limited \\
    Buoy Health & Symptom-checking chatbot & Guided symptom triage & Questionnaire-based; no document analysis & No \\
    Ada Health & AI-powered symptom assessment & Personalised health insights & No structured report parsing & No \\
    WebMD / Healthline & Encyclopaedic content & Comprehensive health reference & Generic content; not personalised to user's report & No \\
    ReportRx (Proposed) & Dual-LLM + RAG (LlamaIndex) & Grounded, structured, disclaimered output & Academic scope; English only & Yes \\
    \hline
\end{tabular}
\end{table}

The comparative analysis reveals that ReportRx occupies a unique position in the landscape: it is the only system in the comparison that combines structured document upload and parsing with AI-powered analysis, RAG-based grounding to minimise hallucination, and a privacy-preserving ephemeral processing model. While general-purpose LLMs offer conversational medical knowledge, they lack the document-centric architecture required for systematic report analysis. Conversely, dedicated medical tools such as Buoy and Ada focus on symptom assessment rather than diagnostic report interpretation.

\section{Feasibility Study}

\subsection{Technical Feasibility}
ReportRx is built entirely on well-documented, widely adopted technologies, each of which has been proven in production at scale. Next.js (version 14) provides a production-grade React framework with server-side rendering, static site generation, and API route capabilities. Its mature ecosystem supports responsive design, component-based development, and seamless deployment to Vercel's edge network. FastAPI is a high-performance Python web framework leveraging asynchronous IO (ASGI) and Pydantic-based data validation, making it suitable for building concurrent REST APIs that can handle multiple simultaneous report analysis requests without blocking. LlamaIndex offers a mature RAG orchestration layer with built-in support for over 40 document loaders, multiple vector store backends (Chroma, Pinecone, Weaviate), and flexible query engine configurations. LLM access is mediated through the Anthropic and OpenAI APIs, both of which provide stable, versioned interfaces with generous rate limits suitable for academic development. All core components are either open-source or accessible via pay-as-you-go API pricing, confirming technical feasibility within an academic setting.

One potential technical risk is the dependency on external LLM API availability. If either the Anthropic or OpenAI API experiences downtime, the system's core analysis pipeline would be unavailable. This risk is partially mitigated by the dual-LLM architecture, which provides two independent LLM backends. Additionally, the use of LlamaIndex's modular design allows the LLM provider to be swapped with minimal code changes, supporting future migration to locally hosted models if required.

\subsection{Economic Feasibility}
As an academic project, core development uses free and open-source tools exclusively. The technology stack --- Next.js, FastAPI, LlamaIndex, PyMuPDF, Tesseract --- carries no licensing costs. LLM API usage during development and testing is cost-manageable; at typical academic usage volumes (approximately 200--400 analyses per month for development and testing), monthly API costs are expected to remain under Rs.\,1,500 ($\sim$\$18 USD). The Next.js frontend can be hosted on Vercel at no cost under its free tier (100 GB bandwidth, 100 hours of serverless function execution per month), and the FastAPI backend can run on a free-tier cloud instance (e.g., Render or Railway) during the project period. No proprietary software licences, hardware purchases, or specialised infrastructure expenditures are required.

If the system were to be scaled beyond academic use, the primary cost driver would be LLM API usage. At production scale (thousands of analyses per month), optimisations such as prompt caching, response caching for similar report patterns, and model distillation could reduce per-query costs by an estimated 40--60\%. Local hosting of a quantised BioMedLM variant could further eliminate the biomedical model API cost without compromising clinical accuracy.

\subsection{Operational Feasibility}
The application interface is designed to require no prior technical knowledge from end users. The single primary interaction --- a drag-and-drop or click-to-upload file field --- is a universally understood web interaction pattern. The structured output format --- covering summary, key observations, out-of-range values, and guidance steps --- is intuitive and self-explanatory. The colour-coded flagging system (red for high values, blue for low values) provides immediate visual comprehension of abnormal results. The mandatory disclaimer at the end of every analysis, rendered in a visually distinct box with a contrasting background, ensures users understand the non-clinical nature of the output.

The system is operationally feasible for the intended audience of non-technical patients for several additional reasons: (i) no account creation or login is required, eliminating a common barrier to adoption; (ii) the application is fully responsive and functions on smartphones, which are the primary internet access device for many target users in the Indian context; (iii) the analysis completes within 30--45 seconds, a duration that is acceptable for a single document analysis task; and (iv) the output can be printed, saved as PDF, or shared directly from the browser using native sharing APIs, enabling users to share the analysis with their healthcare provider.

